% ch06-alternatives.tex – Chapter 6: design alternatives and decision matrix (final round, Writer).
\ChapterMission{המשימה: להשוות שלוש חלופות שונות במהותן ולבחור ביניהן בשקיפות}
\ChapterIcon{info}
\chapter{חלופות תכן ומטריצת החלטה}\label{chap:alternatives}
\begin{scenebox}
על שולחן ישיבות שלושה דפים. על הראשון רמזור עם ספרות שסופרות את האדום, על השני חיישן ובקר, על השלישי כרזה.
מישהו מחק בקו עבה את המילה הראשונה בכותרת של הדף הראשון וכתב מעליה שאלה.
\end{scenebox}

הדף שנמחק מייצג את הרעיון המקורי, והשאלה שמעליו היא נושא הפרק: האם הרעיון עומד מול הראיות. הפרק משווה שלוש חלופות שונות במהותן במטריצה משוקללת ומנמק את הבחירה.

\section{שלוש החלופות}
\paragraph{חלופה \en{A1}: ספירה לאחור בזמן האדום.}
זהו הרעיון המקורי: תצוגה שמראה לממתין כמה זמן נותר עד הירוק, מתוך הנחה שידיעת הזמן תפחית את חוסר הסבלנות.
העדות לטובתה קיימת: בהשוואה בין שני מעברים דומים, במעבר שהייתה בו תצוגת ספירה היו פחות חוצים באדום \cite{lipovac2013jte}, אם כי מן התקציר לא ברור באיזה מופע פעלה התצוגה.
מנגד, התקן מתיר ספירה להולכי רגל רק במופע הפינוי \cite{fhwa2025mutcd}; ליד בתי ספר בסין הלכה ספירה בזמן האדום יחד עם יותר הפרות ויותר ריצה של ילדים \cite{fu2016}; ובשישה-עשר צמתים בסין לא נמצא קשר מובהק בין קיום ספירה לשיעור ההפרות \cite{xing2022}.
ספירה מציגה את ההמתנה ואינה מקצרת אותה, והולכים נטו לחצות מיד כשנותר זמן רב עד הירוק \cite{shah2024}; לפי הסקת הצוות, תצוגה החושפת המתנה ארוכה עלולה אף להאיץ חצייה.

\paragraph{חלופה \en{A2}: אדום מסתגל, ספירה בפינוי וקמפיין.}
חישה בשפת המדרכה סופרת ממתינים ומודדת את זמן המתנתם, והבקר מקצר את האדום בגבולות הבטוחים של פרק~\ref{chap:requirements}, בעדיפות להולך יחיד; הספירה מוצגת רק במופע הפינוי, וקמפיין מציג את ההמתנה כנורמה.
היא פועלת על הגורם שמאמר הדגל זיהה, הולך יחיד וממתין \cite{rosenbloom2009}, ועל הממצא שהציות יורד כשההמתנה מתארכת \cite{vanhouten2007}; חוקרים שמצאו שכל שנייה של אדום מעלה את סיכויי ההפרה המליצו על תזמון שמגיב למספר ההולכים \cite{zhu2021}.

\paragraph{חלופה \en{A3}: קמפיין בלבד.}
קמפיין שליטה חברתית באתר ובתקשורת, בלי חומרה ובלי שינוי בתזמון. מאמר הדגל ממליץ על קמפיין כזה \cite{rosenbloom2009}, אך השפעתו לא נבדקה, והוא אינו מודד את ההמתנה ואינו משנה אותה.

\section{מטריצת ההחלטה}
כל חלופה $j$ מקבלת ציון $c_{ij}$ בסולם \range{1}{5} בכל קריטריון $i$, והציון המשוקלל שלה מחושב לפי משוואה~\eqref{eq:score}:
\begin{equation}
  S_j = \sum_{i} w_i \, c_{ij}, \qquad \sum_{i} w_i = 1
  \label{eq:score}
\end{equation}
המשקלות נגזרים ממטרת הפרויקט: ההשפעה הצפויה על חצייה באדום מקבלת \num{0.30}, משום שזו המטרה; התאמה לתקן מקבלת \num{0.20}, משום שחלופה שלא תאושר לא תופעל; סגירת פער הידע, כלומר מדידת ההמתנה, מקבלת \num{0.20}, משום שזו מגבלת מאמר הדגל. תנועה, עלות ופרטיות מקבלות \num{0.10} כל אחת, משום שהדרישות כבר מגבילות אותן. טבלה~\ref{tab:matrix} מציגה את הציונים ואת נימוקם.

\begin{table}[ht]
  \centering
  \caption{מטריצת החלטה משוקללת (ציונים \range{1}{5}; הציונים הם הערכות הצוות על סמך המקורות שבטקסט)}
  \label{tab:matrix}
  \small
  \begin{tabularx}{\linewidth}{Xcccc}
    \toprule
    \headrow \thead{קריטריון} & \thead{משקל} & \thead{\en{A1}} & \thead{\en{A2}} & \thead{\en{A3}} \\
    \midrule
    השפעה צפויה על חצייה באדום & \num{0.30} & \num{2} & \num{4} & \num{2} \\
    התאמה לתקן ולאישור & \num{0.20} & \num{1} & \num{4} & \num{5} \\
    סגירת פער הידע: מדידת המתנה ונתונים & \num{0.20} & \num{1} & \num{5} & \num{1} \\
    השפעה על תנועת כלי הרכב & \num{0.10} & \num{5} & \num{3} & \num{5} \\
    עלות ומורכבות & \num{0.10} & \num{4} & \num{2} & \num{5} \\
    פרטיות ומקובלות ציבורית & \num{0.10} & \num{5} & \num{3} & \num{5} \\
    \midrule
    ציון משוקלל $S_j$ & \num{1.00} & \num{2.40} & \textbf{\num{3.80}} & \num{3.30} \\
    \bottomrule
  \end{tabularx}
\end{table}

\en{A2} מקבלת את הציון הגבוה ביותר, \num{3.80}, ו-\en{A1} את הנמוך ביותר, \num{2.40}, משום שהיא נכשלת בתקן ובמדידת ההמתנה.
הבחירה יציבה: גם כשהמשקל המשותף של ההשפעה ושל פער הידע יורד מ-\num{0.50} ל-\num{0.20}, \en{A2} נשארת ראשונה (\num{3.68} מול \num{3.36}). רק במשקלות שווים לכל הקריטריונים \en{A3} מובילה (\num{3.83} מול \num{3.50}), אך פירוש הדבר שעלות שקולה לבטיחות.
מכיוון ש-\en{A3} היא חלק מ-\en{A2}, היא גם תוכנית הגיבוי אם האישור לקיצור האדום יתעכב (פרק~\ref{chap:risks}).
כך מאשרת המטריצה את שינוי הכיוון: הפרויקט התחיל מ-\en{A1}, והתקן והראיות, ולא העדפה, הם שהובילו ל-\en{A2}.

